\chapter{Fundamentação Teórica}
\label{cap:FundamentacaoTeorica}

Este capítulo apresenta a fundamentação teórica necessária ao desenvolvimento da pesquisa, abordando os principais conceitos, técnicas e tecnologias relacionados ao tema. Cada seção deve corresponder a um conceito central para a compreensão da metodologia proposta -- evite incluir conceitos que não serão retomados na metodologia ou nos resultados.

\section{Nome do Primeiro Conceito}

Apresente aqui o primeiro conceito teórico relevante para o trabalho. Utilize citações para embasar as definições apresentadas. Há duas formas principais de citar um autor conforme a ABNT:

Na forma narrativa, o nome do autor faz parte da frase: \citeonline{marconi2017metodologia} definem o método científico como um conjunto de procedimentos sistemáticos para a produção de conhecimento válido.

Na forma parentética, a citação aparece entre parênteses ao final da frase, sem fazer parte da sintaxe da oração \cite{marconi2017metodologia}.

Para citar uma página específica de uma obra (recomendado em citações diretas), utiliza-se \citeonline[p.~10]{marconi2017metodologia}.

A Figura \ref{fig:exemplo} ilustra um exemplo de diagrama simples, construído diretamente em \LaTeX{};

\begin{figure}[H]
    \centering
    \caption{Exemplo de diagrama simples}
    \label{fig:exemplo}
    \centering
    \begin{tikzpicture}[
        node distance=2cm, 
        every node/.style={draw, rectangle, minimum width=2.5cm, minimum height=1cm, align=center, rounded corners=3pt},
        >=Stealth
    ]
        \node (a) {Elemento A};
        \node (b) [right=of a] {Elemento B};
        \node (c) [right=of b] {Elemento C};
        
        \draw[->, thick] (a) -- (b);
        \draw[->, thick] (b) -- (c);
    \end{tikzpicture}
    \fonte{Elaborado pelo autor, 2026.}
\end{figure}

\section{Nome do Segundo Conceito}

Apresente aqui o segundo conceito teórico, por exemplo um algoritmo relevante para o trabalho. O Algoritmo \ref{alg:exemplo} ilustra o uso do pacote \textit{algorithm2e} (já configurado em português no preâmbulo) para descrever pseudocódigo.

OBS: Se citar figura ou algoritmo, tem que informar a página do livro ao qual referenciou na fonte da imagem ou algoritmo. ex:

\begin{algorithm}[H]
    \caption{Nome do algoritmo de exemplo}
    \label{alg:exemplo}
    \KwIn{$n$: um valor de entrada}
    \KwOut{$r$: o resultado calculado}
    $r \gets 1$\;
    \For{$i \gets 1$ \KwTo $n$}{
        $r \gets r \times i$\;
    }
    \Return{$r$}\;
\end{algorithm}
\fonte{\citeonline[p.~XX]{marconi2017metodologia}. Substitua pela referência e página reais do livro de onde o algoritmo foi adaptado.}

\section{Nome do Terceiro Conceito}

Se o trabalho envolver algum tipo de código-fonte (consultas, scripts, trechos de configuração), utilize \verb|\lstinputlisting| para importar o código de um arquivo externo, mantendo-o destacado do texto corrido e numerado como Listagem. A Listagem \ref{lst:exemplo} apresenta um exemplo.

\lstinputlisting[
    language=SQL,
    caption={Exemplo de criação de tabela},
    label={lst:exemplo}
]{sql/department.sql}
\vspace{-0.4cm}
\fonte{\citeonline[p.~14]{silberschatz2020database}.}

\section{Trabalhos Relacionados}

Encerre a fundamentação teórica com uma seção que situe o trabalho perante a literatura já existente sobre o tema. Para cada trabalho relacionado relevante, apresente: (i) o que o autor fez; (ii) qual metodologia ou abordagem foi utilizada; (iii) quais os principais resultados; e (iv) como esse trabalho se relaciona com a pesquisa proposta.

Por exemplo: \citeonline{marconi2017metodologia} discutem [aspecto relevante], utilizando [metodologia/abordagem], concluindo que [principal resultado]. Esse achado é relevante para este trabalho pois [conexão direta com a pesquisa proposta].

Ao final da seção, é recomendável indicar explicitamente a lacuna identificada na literatura que este trabalho pretende preencher, deixando claro o ineditismo ou a contribuição da pesquisa.
